% META: {"id":"1SPE-GEOREP-CO-026","chapitre":"1SPE-GEOMETRIE-REPEREE","type_objet":"corrige","exercice_ref":"1SPE-GEOREP-EX-026","capacites_codes":["C3"],"sources_inspiration":[],"status":"approved","origin":"generated","status_history":["generated","audited","accepted","approved"],"release_acceptance":"RELEASE_OWNER_BATCH_ACCEPTANCE","acceptance_closure_digest":"sha256:9b3ccf9a81c5520fbb7e03b7b2d3e7bf2057a02834b6d908bf3be5f49f3b553f","acceptance_receipt_digest":"sha256:4fad86f0282e0fa4e0419cca1ad6379ae7af5a280c04a4a90880f90a1c044242"}
% BEGIN-VERIFY
% from sympy import *
% x, y = symbols('x y')
% # x^2 + y^2 - 6x + 2y + 15 = 0
% # A=-6, B=2, C=15
% # r^2 = 9 + 1 - 15 = -5 < 0 : pas un cercle
% assert 9 + 1 - 15 == -5
% END-VERIFY
\begin{corrige}{1SPE-GEOREP-EX-026}

$A = -6$, $B = 2$, $C = 15$.

$\dfrac{A^2}{4} + \dfrac{B^2}{4} - C = 9 + 1 - 15 = -5 < 0$.

L'équation n'a pas de solution : elle ne représente \textbf{aucun cercle}.

\end{corrige}
